\chapter{Teorija praćenja, pohrane i \mbox{vizualizacije} mrežnog prometa}

Sigurnosni nadzor mreže \engl{Network Security Monitoring - NSM} predstavlja metodu prikupljanja, analize i eskalacije indikatora koji ukazuju na maliciozne aktivnosti unutar nekog mrežnog okruženja. 
Tradicionalni mehanizmi zaštite poput vatrozida i antivirusa se fokusiraju na prevenciju napada izgradnjom barijera tj. statičke obrane. Usprkos tome, kako ističe 
Bejtlich \cite{bejtlich_nsm}, svaka obrana s vremenom zakaže zato što napadači stalno traže nove načine za zaobilaženje tih obrana - vječna igra mačke i miša.
NSM se zato ne oslanja isključivo na spriječavanje i statičku obranu nego i na dubinsku vidljivost mrežnog prometa koja omogućava pravovremenu detekciju i reakciju odnosno aktivnu obranu za incidente.

Kako bi se uspostavio funkcionalan NSM sustav, potrebno je integrirati nekoliko različitih klasa alata: od mrežnih senzora za ekstrakciju metapodataka, preko optimiziranih baza podataka za 
pohranu vremenskih nizova, do grafičkih sučelja za vizualizaciju anomalija. 
U nastavku ovog poglavlja bit će opisane tehnologije koje su odabrane kao temelj za izgradnju takvog rješenja.

\section{Mrežna infrastruktura i pristupne točke}

AP \engl{Access Point} odnosno pristupna točka je uređaj koji se ponaša kao most \engl{bridge} između bežičnih uređaja i usmjernika te time omogućava uređajima poput laptopa, mobitela, IoT senzora 
da se povežu sa žičanom lokalnom mrežom. Za razliku od potrošačkih bežičnih usmjernika, koji u jednom kućištu objedinjuju usmjeravanje, filtriranje prometa i bežično povezivanje, pristupna se točka ograničava 
isključivo na bežično povezivanje.

Alat \texttt{hostapd} omogućuje da se odabrano bežično sučelje na Linuxu pretvori u pristupnu točku \cite{hostapd_docs}. Među njegovim mogućnostima ističu se:
\begin{itemize}
    \item upravljanje pristupnom točkom prema normama serije IEEE 802.11;
    \item podrška autentifikacijskim protokolima poput IEEE 802.1X, WPA, WPA2, WPA3;
    \item rad u ulozi RADIUS poslužitelja ili klijenta, što je značajno korisnicima kojima je potrebna centralizirana autentikacija.
\end{itemize}

Uređaji povezani na pristupnu točku trebaju dobiti mrežne parametre i mogućnost razlučivanja imena domena. Te usluge u sustavima s ograničenim resursima često pruža alat \texttt{dnsmasq} \cite{dnsmasq_docs}. 
Potrebno je naglasiti da \texttt{dnsmasq} ne obavlja usmjeravanje ili filtriranje mrežnog prometa, što su primarni zadatci usmjernika odnosno vatrozida, 
već objedinjuje funkcionalnosti DNS prosljeđivača \engl{DNS forwarder} i DHCP poslužitelja. Alat prihvaća lokalne DNS upite te ih prosljeđuje vanjskim DNS poslužiteljima kako bi se imena 
(domene) prevela u odgovarajuće IP adrese. 

\section{Sustavi za detekciju upada i analizu prometa}

Zeek je platforma otvorenog koda \cite{zeek_docs} namijenjena dubinskoj analizi mrežnog prometa. Za razliku od pristupa koji bilježe sirove pakete (Wireshark), Zeek za svaku promatranu vezu generira strukturirani transakcijski
zapis \engl{transaction log} s metapodatcima poput trajanja veze, broja prenesenih okteta i korištenog protokola. Format izlaznih zapisa moguće je prilagoditi, a analiza se
provodi neovisno o veličini mreže.

Zeek posjeduje vlastiti skriptni jezik kojim se definiraju pravila detekcije i određuje koji se metapodatci izdvajaju iz prometa. Time se logika detekcije prilagođava
specifičnostima nadzirane mreže bez izmjene jezgre sustava.

Zeek razvrstava zapise prema protokolu i vrsti zabilježene aktivnosti (\texttt{conn.log} za sve zamijećene veze, \texttt{dhcp.log} za dodjelu lokalnih IP adresa uređajima, \texttt{http.log} za praćenje HTTP zahtjeva, \texttt{files.log} za praćenja prijenosa datoteka, \texttt{dns.log} za praćenja prevođenja simboličkih imena u adrese i drugi).
Uz zapise vezane uz pojedini protokol, okvir za obavijesti \engl{notice framework} bilježi izvedene sigurnosne događaje u datoteku \texttt{notice.log}.
Detekcijska se logika i popis prijetnji zadaju kroz \texttt{local.zeek} datoteku koja sadrži pravila detekcije napisana Zeekovim skriptnim jezikom, dok \texttt{intel.dat} datoteka sadrži popis indikatora kompromitacije koji se učitava neovisno o toj logici.

Alternativni alati poput Suricate \cite{suricata_docs} primarno su usmjereni na detekciju temeljenu na potpisima \engl{signature-based detection}. Sveobuhvatna rješenja poput Security Oniona zahtijevaju znatno više resursa, najmanje 12 GB radne memorije za minimalnu instalaciju \cite{security_onion, security_onion_hwreq}, zbog čega nisu prikladna za 
potrošačko sklopovlje. Prednost je Zeeka u olakšanoj ekstrakciji metapodataka i svestranom skriptnom jeziku koji omogućuje izradu vlastite detekcijske logike uz malen memorijski otisak.

\section{Kontejnerizacija mrežnih servisa}
Docker je platforma za razvoj, dijeljenje i pokretanje aplikacija. Docker omogućava odvajanje aplikacija od infrastrukture \cite{docker_docs}.

Kontejneri su izolirana okruženja koja se izvode na računalu domaćinu, odvojena od ostalih procesa korištenjem prostorâ imena \engl{namespaces} i kontrolnih grupa \engl{control groups, cgroups}. Svaki je kontejner instanca 
slike \engl{image} koju je moguće izraditi, pokrenuti, zaustaviti, premjestiti ili izbrisati Docker alatima. Slike su prenosive, pa se kontejneri mogu izvoditi lokalno, u virtualnim strojevima i u oblaku.

Zbog izolacije prostorima imena i kontrolnim grupama kontejner ima malen utjecaj na domaćina, a za njegovo pokretanje nije potrebno instalirati dodatne ovisnosti na domaćinu. Svaki kontejner je 
izolirano okruženje te može komunicirati sa ostalim kontejnerima tek ako mu se to eksplicitno dozvoli.

Za orkestraciju više kontejnera koji zajedno čine jednu aplikaciju Docker nudi alat Docker Compose. Konfiguracija cijelog skupa servisa definira se u jednoj 
\textit{docker-compose.yml} datoteci koja opisuje svaki servis, njegove ovisnosti, volumene i mrežne postavke. Pokretanje cijelog sustava svodi se na jednu naredbu 
\texttt{docker compose up -d}, što znatno pojednostavljuje upogonjavanje i reproducibilnost okruženja te je glavni razlog zašto je upravo Docker korišten u ovom radu.

\section{Pohrana vremenskih nizova i vizualizacija}
InfluxDB je baza podataka temeljena na modelu vremenskih nizova \engl{time-series}, optimizirana za pohranu, indeksiranje i upite nad podatcima kod kojih je vrijeme glavna
organizacijska dimenzija \cite{influxdb_docs}. Glavne karakteristike InfluxDB su:
\begin{itemize}
    \item indeksiranje temeljeno na vremenu;
    \item pohrana u kojoj je moguće samo dodavanje \engl{append};
    \item optimizirana za velik broj sekvencijalnih upisa;
    \item uklanjanje podataka provodi se politikama zadržavanja \engl{retention policy}, a ne pojedinačnim izmjenama;
    \item kompresija i smanjenje uzorkovanja \engl{downsampling} radi manjeg zauzeća prostora;
    \item ugrađene agregacijske funkcije nad vremenskim intervalima (pomični prosjeci, percentili, sažetci).
\end{itemize}

U ovom je radu korišten Flux, funkcijski skriptni jezik otvorenog koda za upite, analizu i obradu podataka vremenskih nizova te izvorni upitni jezik za InfluxDB 2.x \cite{flux_getstarted, flux_basics_video}. 
Za razliku od InfluxQL-a, koji je sintaksno blizak SQL-u, Flux se temelji na kompoziciji operacija povezanih u cjevovod te uz InfluxDB podržava i druge izvore podataka, primjerice PostgreSQL i BigQuery, kao i 
formate CSV i JSON.

Primjer Flux upita prikazan je u ispisu \ref{lst:basic_flux_query} \cite{flux_query_basics}.
\begin{lstlisting}[caption={Primjer Flux upita}, label={lst:basic_flux_query}]
import "influxdata/influxdb/v1"

option v = {bucket: "my-bucket", org: "my-org"}

from(bucket: v.bucket)
    |> range(start: -1h)
    |> filter(fn: (r) => r._measurement == "cpu")
    |> mean()
\end{lstlisting}
Upit \ref{lst:basic_flux_query} definira spremnik informacija \engl{bucket} te iz njega redom povlači informacije na sljedeći način: svi zapisi iz prethodnih sat vremena (\texttt{start: -1h}),
koji u sebi sadrže informaciju o CPU-u (\texttt{filter(fn: (r) => r.\_measurement == "cpu")}) te potom iz tih informacija izvlači srednja vrijednost (\texttt{mean()}). Drugim riječima ovaj upit će
prikupiti informacije o CPU-u u zadnjih sat vremena te će prikazati njihovu srednju vrijednost.

Telegraf je agent otvorenog koda za prikupljanje podataka koji razvija InfluxData, a koji prikuplja, obrađuje i agregira metrike, zapise i druge podatke iz različitih izvora. 
Telegraf služi kao primarni sloj za unos podataka u InfluxDB, bazu podataka vremenskih serija, koristeći arhitekturu temeljenu na dodacima s preko 300 ulaznih i izlaznih dodataka za prikupljanje 
podataka iz sustava, senzora i aplikacija te njihov izravan zapis u InfluxDB \cite{telegraf_docs}\cite{influxdb_telegraf_intro}.
Neke od karakteristika Telegrafa su:
\begin{itemize}
    \item arhitektura temeljena na dodacima;
    \item napisan u jeziku Go;
    \item pretprocesiranje podataka;
\end{itemize}

Uz InfluxDB i Telegraf, kao alternativa ispitan je i Elasticsearch s Filebeat agentom za prikupljanje podataka.

Elasticsearch je distribuirani sustav za pretraživanje i analizu podataka otvorenog koda temeljen na Apache Lucene biblioteci \cite{elasticsearch_docs}. Za razliku od 
InfluxDB-a koji optimizira pohranu vremenskih serija, Elasticsearch tretira svaki zapis kao JSON dokument koji se pohranjuje kroz invertirani indeks  
\engl{inverted index} - struktura podataka koja omogućava brzu pretragu po sadržaju polja. Ta arhitektura čini Elasticsearch pogodnim za 
pretraživanje punim tekstom i korelaciju podataka iz heterogenih izvora, no dolazi s većim memorijskim i prostornim otiskom u usporedbi s InfluxDB-om što će biti detaljnije 
prikazano u poglavlju evaluacije.

Kao što Telegraf unosi podatke u InfluxDB, Filebeat preuzima analognu ulogu u ekosustavu Elasticsearcha \cite{filebeat_docs}. Filebeat je lagan agent koji prati zadane datoteke 
zapisnika i nove zapise šalje Elasticsearchu, koji ih potom indeksira. U ovom radu Filebeat čita Zeekove JSON zapise iz direktorija dijeljenog s ostalim servisima i prosljeđuje ih 
Elasticsearchu bez transformacije.

Grafana je platforma otvorenog koda za vizualizaciju i analizu podataka koja podržava više od pedeset izvora podataka, uključujući InfluxDB i 
Elasticsearch \cite{grafana_docs}. Kontrolne ploče se konfiguriraju kroz grafičko sučelje i mogu se izvesti kao JSON datoteke te pohraniti u repozitorij, što 
omogućava automatsko učitavanje pri pokretanju sustava \engl{provisioning}. U ovom radu Grafana služi kao jedinstveno vizualizacijsko sučelje za oba ispitana 
pozadinska rješenja.